\section{Original-field character comparison and the complete PointHV assembly}
\label{sec:point-hv-assembly}

\status{HV1--HV10 prove the entire original-field assembly of
\texttt{PointHV}.  The two auxiliary count suppliers
\texttt{CountVs} and \texttt{CountIV} are proved here; \texttt{CountZ}
is already supplied by CZ6.  The precise remaining count premise is
\texttt{CountI}$\wedge$\texttt{CountII}$\wedge$\texttt{CountIII}
$\wedge$\texttt{CountV}.  No character comparison, layer descent,
case partition, finite-prime summation or uniformity-of-constants
premise remains in this implication.  In particular, the comparison
does not discard primes dividing the original field degree.
The arithmetic theorems that supply the remaining four counts are
separate obligations.}

\subsection{Notation and the exact interfaces}

Let $\chi$ be the primitive Dirichlet character of the source, put
$q=\operatorname{cond}(\chi)$ and $o=\operatorname{ord}(\chi)$,
and let $h$ be a host with exponent $s\ge2$ and prime pair $q_1,d$.
Thus $q_1\equiv3\pmod8$, $d\equiv7\pmod{16}$,
$(q_1d,q)=1$, and
\[
 N=\operatorname{lcm}(2^{s+2},q)q_1d.
 \tag{HV.1}
\]
Inside $(\mathbb Z/N)^\times$ let $H$ be the simultaneous kernel
of the four conditions: reduction modulo $2^{s+2}$ lies in
$\langle-1\rangle$; the $\chi$ value is $1$; and the reductions
modulo $q_1$ and $d$ are squares.  Write $n=[(\mathbb Z/N)^\times:H]$.
The finite exceptional set is the literal source set
\[
 \mathcal B_\chi=\{3\}\ \cup\
 \{\text{odd prime divisors of }o q p_K h_* p_1\},
 \tag{HV.2}
\]
where $p_K=p_K(\chi)$, $h_*=h_*(\chi)$ and $p_1=p_1(\chi)$
are the fixed quantities of \texttt{StmtFields}.
The condition \texttt{Big}$(t)$ says that $q_1,d>2t+1$ and that
each of $q,o,p_K,h_*,p_1$ is smaller than both $q_1$ and $d$.
All these constants have their original source meanings.

For an odd prime $\ell$, let $O$ be any source coefficient DVR
inside the fixed $\ell$-adic algebraic closure, finite over
$\mathbb Z_\ell$, and put
\[
 a_O=\operatorname{length}_O(O/\ell O)>0,
 \qquad A_O=O\otimes_{\mathbb Z_\ell}
                         (\mathbb Q_\ell/\mathbb Z_\ell).
 \tag{HV.3}
\]
For a character $\psi$ modulo $N$ trivial on $H$, the relevant
representation on $A_O$ is
\[
                 \rho_{e,\psi}=\psi^{-1}\chi_{\rm cyc}^{-e}.
 \tag{HV.4}
\]
Let $S'_\ell(\psi,e)$ be its continuous $H^1$ classes over
$\mathbb Q$ which are unramified outside $\ell N$ and locally
zero at every prime dividing $\ell N$.  Its length and normalized
length are exactly
\[
 L_\ell(\psi,e)=\mathtt{Vlen}(\psi,e),\qquad
 v_\ell(\psi,e)=L_\ell(\psi,e)/a_O=\mathtt{Vn}(\psi,e).
 \tag{HV.5}
\]
The quotient notation for the length is used only after finiteness
has been proved; before then we use extended lengths.  Let $X_e$
be the characters trivial on $H$ with $\psi(-1)=(-1)^e$.
The assertion \texttt{VLe}$(\psi,e,B)$ is precisely that
$L_\ell(\psi,e)$ is finite and $v_\ell(\psi,e)\log\ell\le B$.

For a number field $F$ we write $s_\ell(F,e)$ for the strict
negative-twist Selmer length (TW.1).  In particular
\texttt{VVat}$(\ell,F,t,B)$ bounds the sum of these Selmer
lengths for $1\le e<t$; it is not a different normalization
of the compact $H^2$ length.  The layer condition
\texttt{LayerLe}$(\chi,\ell,j,t,B)$ requires this bound on
every model of the $j$th cyclotomic $\mathbb Z_\ell$-layer of
$\mathbb Q(\mu_{\ell q})$.  This is the original universal
quantifier over \texttt{IsAbField} and \texttt{IsLayer}.

The only classical arithmetic inputs below are the elementary
cyclotomic local-field description, the abelian
conductor--discriminant formula already matched in TW-CD, and
the previously proved comparison TW3 and bound KB8.  All
restriction errors and character decompositions are proved here.
For the cyclotomic local description one can use the usual
Eisenstein polynomial $\Phi_{\ell^r}(1+X)$; the valuation argument
needed below is included in HV2.  The cohomological exact
sequence used below is restriction--inflation, with its
transgression; its cochain construction is recalled explicitly.
For the finite cyclotomic and local-field facts see Milne,
\emph{Algebraic Number Theory}, v3.08, Proposition 6.2 and
Proposition 7.55,
\url{https://www.jmilne.org/math/CourseNotes/ANT.pdf}.
The conductor--discriminant formula is equation (4.2.1) in
Sharifi, \emph{Iwasawa Theory}, Chapter 4,
\url{https://math.ucla.edu/~sharifi/notes/iwasawa-ch04.html}.

\paragraph{HV1: the original host field and its elementary character bounds.}
\label{claim:HV1}
Every original field $E$ satisfying \texttt{IsChiField} for $h$
has, in any algebraic closure of $\mathbb Q$, image
$\mathbb Q(\mu_N)^H$.  In particular it is abelian,
unramified outside $N$, and has degree $n$, with
\[
 n\mid2^{s+2}o,\qquad n\le N,\qquad16\mid N,\qquad
 1\le o\le q,\qquad |X_e|\le N.
 \tag{HV.6}
\]
Write $\psi_p$ for the CRT component of $\psi$ supported at $p$.
For odd $p$, $\operatorname{ord}(\psi_p)\mid2o$; if $p\nmid q$
then $\psi_p^2=1$.  If
$c_2=v_2(\operatorname{cond}(\psi_2))$, then
\[
                   \operatorname{ord}(\psi)\le2o\,2^{c_2}.
 \tag{HV.7}
\]
Write $\psi=\kappa\theta$ for its $\ell$-power and prime-to-$\ell$
order parts, and put
$\chi_P=\theta/\theta_\ell$ and $\theta_1=\theta/\theta_2$,
as in the source.  If $\ell$ is odd, $\ell\ne q_1,d$ and
$\operatorname{ord}(\chi_P)$ is odd, then
\[
          \psi_2=1,\qquad\operatorname{cond}(\psi)\mid q.
 \tag{HV.8}
\]
The same conclusions hold if $\operatorname{ord}(\psi)$ is odd.
Consequently the number of host characters satisfying (HV.8)
is at most $q$, in any coefficient ring.
Under \texttt{Big}$(t)$, every $\ell\in\mathcal B_\chi$ is
smaller than $q_1,d$ and hence smaller than $N$.  If $\ell$ is
odd and outside $\mathcal B_\chi$, then
\[
 \ell\ge5,\quad\ell\nmid oqn;
 \quad 1\le e<t, \ell-1\mid2e
       \ \Longrightarrow\ e\ge2,\ \ell\le2t,\ \ell\nmid N.
 \tag{HV.9}
\]

\begin{proof}
The defining embedding of \texttt{IsChiField} specifies exactly
the automorphisms which fix its image: the sign condition at
$2^{s+2}$, the kernel of $\chi$, and the two quadratic-residue
conditions.  For a unit modulo an odd prime, being a square is
equivalent to Legendre symbol $1$.  If the initially supplied
cyclotomic level is $N'$, its divisibilities give $N\mid N'$.
Its fixing group is the preimage of $H$ under reduction modulo
$N$, since all four defining conditions depend only on those
residues.  In particular it contains the kernel of the restriction
to $\mathbb Q(\mu_N)$, so the image of $E$ already lies in
that smaller field.  Cyclotomic restriction is surjective;
its fixing group there is exactly $H$.
The Galois correspondence gives the asserted field, with no
change to the original $E$.
Different embeddings are conjugate; since the cyclotomic Galois
group is abelian they have this same image.

The quotient by $H$ injects into the product of the quotient
of $(\mathbb Z/2^{s+2})^\times$ by $\langle-1\rangle$, the
image of $\chi$, and two groups of order two.  These have orders
$2^s,o,2,2$, proving the divisibility in (HV.6).
Also $n\le\varphi(N)\le N$, and $o\mid\varphi(q)$ gives
$o\le q$.  A finite abelian group of order $n$ has at most $n$
characters in a characteristic-zero domain: decompose it into
cyclic groups and use the bound of $r$ roots for $X^r-1$.
This proves the bound on $X_e$ and also on the union $X_0\cup X_1$.

On an odd CRT factor the first quotient is trivial, so its
exponent divides $2o$.  If $p\nmid q$, only the quadratic factors
remain.  The component $\psi_2$ factors through the units modulo
$2^{c_2}$, and its order divides $2^{c_2}$.  The product of all
odd components has order dividing $2o$, which proves (HV.7).
There are no components outside the prime divisors of $N$.

For (HV.8), the $2$-component has $2$-power order.  Since $\ell$
is odd, it is unchanged on removing the $\ell$-power order part,
and also on removing the component at $\ell$.  Its order
therefore divides the odd order of $\chi_P$, and it is trivial.
At $q_1,d$, which differ from $\ell$, the component of $\psi$
is a product of an odd-order component of $\chi_P$ and an
$\ell$-power order component of $\kappa$.  It has odd order;
it also has order dividing two, and is trivial.  At every other
odd prime the exponent of that prime in $N$ is exactly its
exponent in $q$.  This proves the conductor divisibility.
For an odd-order $\psi$ the same argument applies directly.
A character of conductor dividing $q$ factors uniquely modulo
$q$, giving at most $\varphi(q)\le q$ possibilities.

All factors defining (HV.2) are positive.  Every prime divisor
is no larger than one of these factors; \texttt{Big} therefore
bounds it by $q_1,d$.  The inserted prime $3$ also satisfies the
bound.  Outside this set the assertions about $o,q,n$ follow
from (HV.6), and $\ell\ge5$.  If $\ell-1\mid2e$, then
$\ell-1\le2e$, so $\ell\le2t$ and $e\ge2$.
The only possible prime divisors of $N$ are $2$, primes of $q$,
$q_1$ and $d$.  The size condition excludes the last two and
the earlier conclusions exclude the first two.  This proves
(HV.9).
\end{proof}

\paragraph{HV2: quantitative restriction without inverting the degree.}
\label{claim:HV2}
Let $F/\mathbb Q$ be finite Galois of degree $b$, and suppose a
finite-order character $\psi$ is trivial on $G_F$.
For $e\ge1$ and odd $\ell$, the restriction kernel on global
$H^1(\mathbb Q,A_O\rho_{e,\psi})$ has normalized length at most
$3b^2e$.  Whenever restriction carries a Selmer group $S$ into
the strict negative-twist group over $F$, one consequently has
\[
       \frac{\operatorname{length}_O S}{a_O}
                         \le s_\ell(F,e)+3b^2e.
 \tag{HV.10}
\]
Finiteness on the right implies finiteness on the left.
The assertion applies to the original algebraic closure and to
every original choice of local places.

\begin{proof}
For a finite extension $K/\mathbb Q_\ell$ of degree at most $b$,
the image $J$ of its cyclotomic character has finite index
$a\le b$ in $\mathbb Z_\ell^\times$.  Indeed the cyclotomic
character of $\mathbb Q_\ell$ is surjective, and the map from
the cosets of $G_K$ to the cosets of $J$ is surjective.
Surjectivity at the base follows from the Eisenstein
polynomials $\Phi_{\ell^r}(1+X)$, which give degree
$\varphi(\ell^r)$ for the local cyclotomic extensions.
The full image is compact and hence closed; its surjectivity
at every finite level therefore gives surjectivity onto
$\mathbb Z_\ell^\times$.
Since the quotient by $J$ has order $a$, $(1+\ell)^a\in J$.
The invariants of $A_O(-e)$ over $K$ are therefore killed by
$(1+\ell)^{ae}-1$.  The binomial formula, successively applied
to an exponent prime to $\ell$ and to powers of $\ell$, gives
\[
 \frac{\operatorname{length}_O A_O(-e)^{G_K}}{a_O}
 \le v_\ell((1+\ell)^{ae}-1)
 =1+v_\ell(a)+v_\ell(e)\le3be.
 \tag{HV.11}
\]
Here the $u$-torsion of $A_O$ has length $a_Ov_\ell(u)$
for $0\ne u\in\mathbb Z_\ell$.  Global invariants over $F$
are a submodule of these invariants at any one place above
$\ell$, so the same bound holds for
$W=A_O(-e)^{G_F}$.

Put $G=\operatorname{Gal}(F/\mathbb Q)$.
The restriction--inflation sequence identifies the restriction
kernel with $H^1(G,W)$, where the quotient action on $W$ is
the one inherited from $\rho_{e,\psi}$, including its
$\psi^{-1}$ factor.  Its underlying $O$-module is the invariant
module just bounded.  In inhomogeneous cochains the degree-one
group consists of $b$ copies of $W$, so
$\operatorname{length}H^1(G,W)\le b\operatorname{length}W$.
This proves the kernel bound, even when $\ell\mid b$.
Restriction of a cocycle preserves each local coboundary and
unramified-inertia condition.  Tensoring the strict group over
$F$ with $O$ multiplies its $\mathbb Z_\ell$-length by $a_O$:
the ring is finite free, so cochains, coboundaries and all local
kernels commute with this tensor product.  These facts prove
(HV.10), including finiteness.

Conjugating a selected place or transporting an algebraic closure
conjugates the cochains and their coboundary witnesses.  The
restriction diagram commutes with this operation.  Thus this
proof concerns the given $F$ and given local data, rather than
only one convenient model.
\end{proof}

\paragraph{HV3: integral decomposition into characters.}
\label{claim:HV3}
Let $G$ be finite abelian, let $O$ contain its character values,
and let $M$ be any $O[G]$-module, initially of possibly infinite
length.  For $c\in G$ of order two, odd residue characteristic,
and $\epsilon=\pm1$, put
$M_\psi=\{x:gx=\psi(g)x\text{ for every }g\in G\}$.
Then
\[
 \operatorname{length}_O M^{c=\epsilon}
 \le\sum_{\psi(c)=\epsilon}\operatorname{length}_O M_\psi.
 \tag{HV.12}
\]
In particular, finiteness of the displayed eigenspaces proves
finiteness of the left side.  No condition on divisibility of
$|G|$ by $\ell$ is imposed.

\begin{proof}
For commuting endomorphisms $f,g$, the map
$g:\ker(fg)\to\ker f$ has kernel $\ker g$.
Additivity and monotonicity of length give
$\operatorname{length}\ker(fg)\le
\operatorname{length}\ker f+\operatorname{length}\ker g$.
If $T^r=1$, factor
$T^r-1=\prod_{\zeta^r=1}(T-\zeta)$ and iterate this inequality.
Apply it successively to cyclic generators in a decomposition
$G=\prod_i C_{r_i}$, starting on $M^{c=\epsilon}$.
All the kernels are stable under the remaining generators.
Every final intersection is one of the simultaneous
eigenspaces $M_\psi\cap M^{c=\epsilon}$.
If $\psi(c)\ne\epsilon$, their difference is the unit $\pm2$,
so that intersection is zero.  This proves (HV.12).
The same successive kernel exact sequences prove finiteness
when all final terms are finite; no arithmetic with infinite
lengths is needed to infer it.
\end{proof}

\paragraph{HV4: the comparison on every original CM field.}
\label{claim:HV4}
Let $E$ be an original CM field of host $h$, and let $D$ be
any of its original local-data records at an odd prime $\ell$.
Fix $e\ge1$.
Choose a source coefficient ring $O$ containing the values of
all characters of $G=\operatorname{Gal}(E/\mathbb Q)$.
If all $L_\ell(\psi,e)$ for $\psi\in X_e$ are finite, then
$\mathtt{plusLen}_D(e)$ is finite and
\[
 \mathtt{plusLen}_D(e)
 \le\sum_{\psi\in X_e}
                 \bigl(v_\ell(\psi,e)+(Ne)^6\bigr).
 \tag{HV.13}
\]
If $\ell\nmid n$, the error can be omitted:
\[
             \mathtt{plusLen}_D(e)
                     \le\sum_{\psi\in X_e}v_\ell(\psi,e).
 \tag{HV.14}
\]
The source's normalization, the original field, and all its
local conditions are retained.

\begin{proof}
Use HV1 to embed the actual $E$ into the fixed rational
algebraic closure.  Let $B$ be its strict Selmer module with
coefficient $A_O(-e)$.  Normality gives the natural $G$-action
on $B$.  If $c$ is complex conjugation, the defining plus
condition, after this identification, is the eigenspace
$B^{c=(-1)^e}$.  To check the sign, the natural action on a
cocycle $z$ is $g\mapsto c\,z(c^{-1}gc)$, and $c$ acts on
the coefficient as $(-1)^e$.  Thus the untwisted conjugation
invariance defining the plus part is precisely the stated
eigenvalue.  Finite free coefficient extension and conjugacy
of the closures give the exact equality
\[
 a_O\mathtt{plusLen}_D(e)
                    =\operatorname{length}_O B^{c=(-1)^e}.
 \tag{HV.15}
\]
This is also the explicit \texttt{partModel} identification.
Changing a local place changes the cocycle and its local
coboundary witness by conjugacy; hence it changes no kernel
in this equality.
For the precise quotient used in $D$, a pure negative-twist
cocycle unramified away from $\ell$ vanishes on those inertia
groups: their action on the coefficient is trivial, so their
coboundaries are zero.  The cocycle equation then makes it
vanish on their conjugates and on the closed normal subgroup
they generate.  It therefore descends to the actual
$G_{E,S_\ell}$ quotient.  Inflation is its inverse on these
classes and preserves every selected local kernel.  This
identifies the absolute-cochain description with the one in $D$.

For a character $\psi$ of $G$, twist the coefficient action by
$\psi^{-1}$.  Over $E$ this does not change the coefficient,
but on $H^1(E,A_O(-e))$ it multiplies the $G$-action by
$\psi^{-1}$.  Thus $B_\psi$ becomes an invariant submodule.
Restriction--inflation gives an exact sequence
\[
 H^1(\mathbb Q,A_O\rho_{e,\psi})
 \xrightarrow{\mathrm{res}}H^1(E,A_O(-e))^G
 \longrightarrow H^2(G,W_\psi).
 \tag{HV.16}
\]
Here the action on the middle term is the twisted one just
specified, and $W_\psi$ has underlying module
$A_O(-e)^{G_E}$ with the quotient action inherited from
$\rho_{e,\psi}$.  In particular this action retains the
$\psi^{-1}$ factor; it need not be the untwisted action.
The sequence can be constructed directly: choose representatives
of $G_{\mathbb Q}/G_E$, express invariance of a cocycle by
coboundary witnesses on these representatives, and compare the
witness for a product with the two individual witnesses.
The resulting degree-two cocycle takes values in
$A_O(-e)^{G_E}$.  Altering the witnesses changes it by a
degree-two coboundary, and its being zero is exactly the
condition that the cocycle extend to $G_{\mathbb Q}$.
All quotient choices are finite, so this is the continuous
restriction--inflation sequence as well.

Let $P$ be the preimage of $B_\psi$ under restriction in
(HV.16).  At a prime $p\ne\ell$, the pure negative twist is
unramified over $E_w$, and
\[
 H^1_{\rm unr}(E_w,A_O(-e))
       =A_O/(p^{-fe}-1)A_O=0,
 \tag{HV.17}
\]
where $f$ is the residue degree.  The scalar is nonzero for
$e\ge1$ and acts surjectively on the divisible module $A_O$.
Thus every strict Selmer class in $B_\psi$ is locally zero
at every finite place of $E$.
If $p\nmid\ell N$, the extension $E_w/\mathbb Q_p$ is
unramified and its inertia is the same subgroup as rational
inertia.  Any lift in $P$ is consequently unramified there.
For $v\mid\ell N$, the localization of such a lift lies in
the restriction kernel
\[
 K_v=H^1(G_v,A_O(-e)^{G_{E_w}}),
       \qquad G_v=\operatorname{Gal}(E_w/\mathbb Q_v).
 \tag{HV.18}
\]
Here too the quotient action on the displayed invariant
module is the one inherited from $\rho_{e,\psi}$.
The kernel of $P\to\prod_{v\mid\ell N}K_v$ is exactly
$S'_\ell(\psi,e)$.  Conversely every class in that rational
Selmer group restricts into $B_\psi$, since local zero at
primes of $N$ gives the unramified condition over $E$ there.
This proves the equality of kernels, including every original
local condition.

The global invariant module in (HV.16) has normalized length
at most $3ne$ by (HV.11).  Degree-two inhomogeneous cochains
have $n^2$ copies, so the cokernel contribution is at most
$3n^3e$.  At $\ell$, (HV.11) and degree-one cochains bound
(HV.18) by $3n^2e$.  At a prime $p\mid N$, $p\ne\ell$,
the normalized invariant length is
\[
 v_\ell(p^{fe}-1)\le fe\log_\ell p\le neN,
 \tag{HV.19}
\]
since $f\le n$, $p\le N$ and $\ell\ge3$.
The corresponding degree-one cochains contribute at most
$n^2eN$.  There are at most $N$ such primes.  We have proved
\[
 \frac{\operatorname{length}_O B_\psi}{a_O}
 \le v_\ell(\psi,e)+3n^3e+3n^2e+Nn^2eN
 \le v_\ell(\psi,e)+7N^4e
 \le v_\ell(\psi,e)+(Ne)^6.
 \tag{HV.20}
\]
Finiteness follows before taking these ordinary lengths: the
local kernels and the degree-two obstruction are finite, and
$S'_\ell(\psi,e)$ is finite by hypothesis.  The last inequality
uses $N\ge16$ and $e\ge1$.
If $\ell\nmid n$, positive-degree cohomology of $G$ and all
$G_v$ is zero on these $\ell$-primary modules.  The averaging
homotopy divides by their orders, which are units in $O$;
equivalently restriction followed by transfer gives the same
vanishing.  Thus every error in (HV.20) is zero in this case.
Combine (HV.12), (HV.15) and (HV.20) to obtain (HV.13)--(HV.14).

There always exists such a coefficient ring: adjoin the finitely
many roots of unity of order dividing $n$ inside the specified
$\ell$-adic algebraic closure and take its ring of integers.
If an initial coefficient ring is prescribed, use the compositum.
For a finite extension $O\subset O'$ of DVRs the length and
$a_O$ are both multiplied by the relative ramification index.
Continuous cochains and all local kernels commute with this
finite free extension, so normalized lengths and their
finiteness are unchanged.  This also verifies compatibility
with the source's unrestricted coefficient quantifiers.
\end{proof}

\subsection{The two auxiliary counts are complete suppliers}

\paragraph{HV5: the complete small-conductor count \texttt{CountVs}.}
\label{claim:HV5}
The literal assertion \texttt{CountVs} holds.  Its exponent can
be chosen as one natural number depending only on $\chi$, before
the host, $t$, prime, coefficient ring, character and algebraic
closure are chosen.

\begin{proof}
In this case $\ell\in\mathcal B_\chi$, $e=1$, and
\[
 2^{c_2}<R_1(\ell,\psi)
      =16o(\ell q)^2 4^{n_0},\qquad
 n_0=v_2(\ell^{f_0}-1),\quad
 f_0=\operatorname{ord}_{\operatorname{lcm}(4,
                      \operatorname{ord}(\theta_1))}(\ell).
 \tag{HV.21}
\]
The order of $\theta_1$ divides $2o$ by HV1 and is prime to
$\ell$, by construction of the prime-to-$\ell$ part.
Consequently the modulus $m=\operatorname{lcm}(4,
\operatorname{ord}(\theta_1))$ divides $4o$, and $\ell$ is
a unit modulo $m$.  Thus $f_0$ is positive,
$f_0\mid\varphi(m)$ and $\varphi(m)\mid\varphi(4o)$.
The last divisibility follows directly from the prime-power
formula for $\varphi$ whenever $m\mid4o$.
It follows that
\[
 2^{n_0}\le\ell^{f_0}-1<\ell^{\varphi(4o)},\qquad
 R_1\le16o(\ell q)^2\ell^{2\varphi(4o)}.
 \tag{HV.22}
\]
Let $M_\chi=\max\mathcal B_\chi$; this is defined since
$3\in\mathcal B_\chi$.  Choose the fixed integer
\[
 b_* = 2o\,[16o(M_\chi q)^2M_\chi^{2\varphi(4o)}]M_\chi.
 \tag{HV.23}
\]
Equations (HV.7), (HV.21) and (HV.22) give
$\ell\operatorname{ord}(\psi)\le b_*$.

Inside the given algebraic closure, take the abelian field
fixed by the intersection of the kernel of $\psi$ and the
kernel of the mod-$\ell$ cyclotomic character.  Denote it by
$F$.  It contains $\mu_\ell$, has degree
$b\le(\ell-1)\operatorname{ord}(\psi)\le b_*$, and lies
in $\mathbb Q(\mu_{\ell N})$.
Restriction carries $S'_\ell(\psi,1)$ into the strict group
over $F$.  Indeed at primes dividing $N$ the rational class
is already locally zero, and at every other prime away from
$\ell$ it is unramified; these properties persist on restriction.
HV2 therefore gives
\[
                  v_\ell(\psi,1)\le s_\ell(F,1)+3b_*^2.
 \tag{HV.24}
\]
The field is totally complex, since $\ell$ is odd and
$\mu_\ell\subset F$.  Apply TW3's strict Selmer comparison
and the full original-field theorem KB8 to this $F$.
With $\Delta_F=|\operatorname{disc}(F)|$ they give finiteness
and
\[
 s_\ell(F,1)\log\ell
 \le\log\bigl|H^2(G_{F,S_\ell},\mathbb Z_\ell(2))\bigr|
 \le\exp\bigl((2b)^{31}\log(2\Delta_F)\bigr).
 \tag{HV.25}
\]
The named abelian conductor--discriminant formula gives
$\Delta_F\le(\ell N)^b$.  Put $x=\ell N\ge2$ and
$A=(2b_*)^{31}(b_*+1)$.  Then
$\log(2\Delta_F)\le(b_*+1)\log x$, so (HV.24)--(HV.25)
imply
\[
             v_\ell(\psi,1)\log\ell
                    \le x^A+3b_*^2\log\ell
                    \le x^A+3b_*^2x.
 \tag{HV.26}
\]
For every integer $b_*\ge1$ one has
$3b_*^2\le2^{2b_*+2}$; this follows, for example, from
$b_*\le2^{b_*}$.
Thus the last expression is at most $x^{A+2b_*+4}$.
Taking
\[
                     C_{\rm Vs}=A+2b_*+4
 \tag{HV.27}
\]
proves the source's \texttt{VLe}$(\ell Ne)^{C_{\rm Vs}}$
in every small-conductor case.  The proof established the
finiteness part first.  Every quantity in (HV.23), (HV.27)
depends only on $\chi$.  This proves the literal \texttt{CountVs},
with a permissible larger exponent than the source construction.
\end{proof}

\paragraph{HV6: the complete layer comparison \texttt{CountIV}.}
\label{claim:HV6}
The literal assertion \texttt{CountIV} holds.  For fixed
$\chi$, $\ell\in\mathcal B_\chi$ and layer index $j$, set
$b_0=\ell q\ell^j$.  One may take
$C_4=2b_0+4$, before the host and all its remaining data.
In detail, for every $t\ge2$ with \texttt{Big}$(t)$ and every
real $B$ satisfying \texttt{LayerLe}$(\chi,\ell,j,t,B)$,
every host character with odd $\operatorname{ord}(\chi_P)$
has finite $L_\ell(\psi,e)$ for $1\le e<t$ and
\[
          \sum_{1\le e<t}v_\ell(\psi,e)
                                  \le B+(Nt)^{C_4}.
 \tag{HV.28}
\]
No parity condition on this character is required in this claim.

\begin{proof}
HV1 gives $\ell\ne q_1,d$ and
$\operatorname{cond}(\psi)\mid q$.
Construct $F_0=\mathbb Q(\mu_{\ell q})$ in the given closure.
It contains $\mu_\ell$ and has degree at most $\ell q$.
Its cyclotomic image is an open subgroup of
$1+\ell\mathbb Z_\ell$, so its root-adjoining tower is a
$\mathbb Z_\ell$-extension.  More explicitly, the valuation
formula from HV2 identifies this group with $\mathbb Z_\ell$;
its successive open subgroups have indices $\ell^j$.
Thus its $j$th layer $F_1$ is a root-adjoining field of exactly
that degree over $F_0$, and
$[F_1:\mathbb Q]\le b_0$.
This constructs an actual field satisfying the source's
\texttt{IsAbField} and \texttt{IsLayer} in the original closure.
The universal layer premise applies to it and gives
\[
       \sum_{1\le e<t}s_\ell(F_1,e)\le B,
       \qquad s_\ell(F_1,e)<\infty.
 \tag{HV.29}
\]
Since the conductor of $\psi$ divides $q$, it is trivial on
$G_{F_1}$.  The same local-condition argument as in HV5 shows
that rational restriction maps $S'_\ell(\psi,e)$ into the
strict Selmer group of $F_1$.  Applying HV2 and summing gives
\[
 \sum_{1\le e<t}v_\ell(\psi,e)
 \le B+3b_0^2\sum_{1\le e<t}e
 \le B+3b_0^2t^2.
 \tag{HV.30}
\]
All these groups are finite by the same restriction argument.
For $N\ge16$, $t\ge2$ and $b_0\ge1$,
$3b_0^2t^2\le(Nt)^{2b_0+4}$; the elementary estimate
$b_0\le2^{b_0}$ proves it immediately.  This is (HV.28).
We used the premise on an actual layer in the original closure;
any supplied layer model can be transported to it by its
$F_0$-isomorphism, which conjugates all cocycles and local
witnesses.  Thus no field-model restriction has entered the
conclusion.  This proves the full \texttt{CountIV}.
\end{proof}

\subsection{The complete partition and summation}

\paragraph{HV7: an exhaustive partition with the original predicates.}
\label{claim:HV7}
Let $\ell$ be odd, $1\le e<t$, assume \texttt{Big}$(t)$,
and let $\psi\in X_e$.  The source cases are exhaustive:
\[
\begin{array}{ll}
\mathrm{I}:&\ell\notin\mathcal B_\chi,\quad\ell-1\nmid2e;\\
\mathrm{Z}:&\ell\notin\mathcal B_\chi,\quad\ell-1\mid2e,
                                                        \quad\psi=1;\\
\mathrm{II}:&e\ge2,\quad\operatorname{ord}(\chi_P)\text{ even},
                 \quad(\ell\in\mathcal B_\chi\text{ or }\ell-1\mid2e);\\
\mathrm{III}:&\ell\notin\mathcal B_\chi,\quad\ell-1\mid2e,
                 \quad\operatorname{ord}(\psi)\text{ odd},\quad\psi\ne1;\\
\mathrm{IV}:&\ell\in\mathcal B_\chi,\quad e\ge2,
                                    \quad\operatorname{ord}(\chi_P)\text{ odd};\\
\mathrm{V}:&\ell\in\mathcal B_\chi,\quad e=1.
\end{array}
\tag{HV.31}
\]
Case V splits into the two disjoint inequalities
$R_1\le2^{c_2}$ and $2^{c_2}<R_1$, to which
\texttt{CountV} and the proved \texttt{CountVs} apply.

\begin{proof}
Inside $\mathcal B_\chi$, first separate $e=1$; otherwise
$e\ge2$ and the positive order of $\chi_P$ is even or odd.
Outside $\mathcal B_\chi$, failure of divisibility is case I.
When divisibility holds, HV1 gives $e\ge2$, $\ell\nmid N$
and $\ell\nmid\operatorname{ord}(\psi)$: the latter follows
from $\operatorname{ord}(\psi)\mid n$ and (HV.6).
Hence $\psi$ already has prime-to-$\ell$ order, and its
$\ell$-component is trivial.  Thus $\chi_P=\psi$.
Separate the trivial character, and then even and odd order,
to obtain Z, II and III respectively.  The numerical split in
case V is the ordinary ordering of the two integers.
\end{proof}

\paragraph{HV8: every exceptional-prime sum, with the prescribed layer.}
\label{claim:HV8}
Assume the actual source statements \texttt{CountII} and
\texttt{CountV}.  Then the literal statement \texttt{SumIn}
holds.  Its constant for $(\chi,\ell,j)$ is chosen before
the host, the twist cutoff, the layer budget, the original
field and all its local data.

\begin{proof}
Fix $\chi$, $\ell\in\mathcal B_\chi$ and $j$.
Take the exponents $C_2,C_5$ provided by the two assumed
counts, $C_{\rm Vs}$ from HV5, and $C_4$ from HV6.
Put
\[
 C_m=\max(C_2,C_5,C_{\rm Vs}),\qquad
 C_{\rm in}=C_4+2C_m+12+\ell q.
 \tag{HV.32}
\]
Now take any subsequent host, $t\ge2$, $B\ge0$ and its
prescribed layer bound, and any original $E,\Omega,D$.
Put $M=Nt\ge32$ and choose a coefficient ring splitting
$\operatorname{Gal}(E/\mathbb Q)$, as in HV4.

Let $Y$ be all host characters for which
$\operatorname{ord}(\chi_P)$ is odd.  By HV1, $|Y|\le q$;
this is a single finite set independent of $e$.
For each $\psi\in Y$, HV6 gives
$\sum_{1\le e<t}v_\ell(\psi,e)\le B+M^{C_4}$.
For every $\psi\in X_e\setminus Y$, HV7 gives case II
when $e\ge2$, or one of the two parts of case V when $e=1$.
The assumed counts and HV5 therefore give
\[
          v_\ell(\psi,e)\log\ell
                \le(\ell Ne)^{C_m}\le M^{2C_m},
 \tag{HV.33}
\]
because $\ell<N$ and $e<t$.
Every character length required by HV4 is finite: use HV6
for $Y$ and these pointwise counts for its complement.
Using $|X_e|\le N$, nonnegativity, and interchanging only
finite sums yields
\[
 \log\ell\sum_{1\le e<t}\sum_{\psi\in X_e}v_\ell(\psi,e)
 \le q\log\ell\,(B+M^{C_4})+M^{2C_m+1}.
 \tag{HV.34}
\]
The total error in (HV.13) is bounded by
\[
 \log\ell\sum_{1\le e<t}|X_e|(Ne)^6
             \le\ell Nt\,M^6\le M^8.
 \tag{HV.35}
\]
Here $\log\ell\le\ell\le N$, and $q\le N$.
Combining (HV.13), (HV.34) and (HV.35) gives
\[
 \log\ell\sum_{1\le e<t}\mathtt{plusLen}_D(e)
 \le\ell q B+M^{C_4+2}+M^{2C_m+1}+M^8
 \le M^{C_{\rm in}}+C_{\rm in}B.
 \tag{HV.36}
\]
For the last inequality each of the three exponents is at
most $C_4+2C_m+8$, and $3M^a\le M^{a+2}$ for $M\ge2$;
also $C_{\rm in}\ge\ell q$.
All plus lengths are finite by HV4.
This is exactly \texttt{PlusLe} for the singleton $\{\ell\}$
and proves \texttt{SumIn}, with every quantifier in its
original order.
\end{proof}

\paragraph{HV9: the sum over all nonexceptional odd primes.}
\label{claim:HV9}
Assume the actual source statements \texttt{CountI},
\texttt{CountII} and \texttt{CountIII}.  Then the literal
statement \texttt{SumOut} holds.  Only the two auxiliary
primes supplied by \texttt{CountI} occur in its host threshold.

\begin{proof}
Fix $\chi$.  Choose $C_1,r_1,r_2$ from \texttt{CountI},
with $r_1,r_2$ prime, and choose the exponents $C_2,C_3$ of
the two other counts.  Put $C_Z=0$ by CZ6 and
\[
       K=\max(C_Z,C_2,C_3),\qquad
       C_{\rm out}=C_1+3K+4.
 \tag{HV.37}
\]
All choices precede the host and field.
Now take a permissible host above $r_1,r_2$, a cutoff $t\ge2$
with \texttt{Big}$(t)$, any original CM field and local data,
and an arbitrary finite set $S$ of odd primes outside
$\mathcal B_\chi$.  Set $M=Nt$.
Choose, for each prime, a coefficient ring splitting the
finite Galois group of this original field.

For a fixed $1\le e<t$, split $S=S_0\sqcup S_1$ according
to whether $\ell-1\nmid2e$ or $\ell-1\mid2e$.
On $S_0$, \texttt{CountI} applies to the whole finite set
and gives the joint bound $M^{C_1}$, as well as finiteness
of every required character length.
For $\ell\in S_1$, HV1 gives $\ell\le2t$.
The cases in HV7 are now Z, II or III; CZ6 and the two
assumed counts give, for every $\psi\in X_e$,
\[
                 v_\ell(\psi,e)\log\ell
                        \le(\ell Ne)^K\le M^{3K}.
 \tag{HV.38}
\]
There are at most $2t$ such primes and at most $N$ characters
at each prime.  Thus the combined character sum at this
weight is at most
\[
                         M^{C_1}+2M^{3K+1}.
 \tag{HV.39}
\]
Every prime of $S$ is coprime to $n$ by HV1.  Consequently
the error-free comparison (HV.14) applies, including at the
possibly very large primes of $S_0$; no sum of an uncontrolled
error over those primes occurs.
Summing (HV.39) over $e$ gives
\[
 \sum_{\ell\in S}\log\ell\sum_{1\le e<t}
                         \mathtt{plusLen}_{D_\ell}(e)
 \le t\bigl(M^{C_1}+2M^{3K+1}\bigr)
 \le M^{C_1+3K+4}.
 \tag{HV.40}
\]
For the last bound use $t\le M$, $M\ge2$, and bound both
terms by the larger displayed monomial, with one extra
power to absorb their sum.
Finiteness follows at every $(\ell,e)$ from the relevant
count and HV4.  Since $S$ was an arbitrary finite set, this
is precisely \texttt{SumOut}, with the original uniform
prime-set quantifier.
\end{proof}

\paragraph{HV10: literal PointHV and immediate window propagation.}
\label{claim:HV10}
The completed suppliers above prove
\[
 \boxed{\quad
 \mathtt{CountI}\wedge\mathtt{CountII}\wedge
 \mathtt{CountIII}\wedge\mathtt{CountV}
                    \ \Longrightarrow\ \mathtt{PointHV}.
 \quad}
 \tag{HV.41}
\]
More explicitly, for each $\chi$ and each prescribed function
$j:\mathbb N\to\mathbb N$, there are natural numbers
$C,r_1,r_2$, with $r_1,r_2$ prime, before any host is chosen,
such that for every permissible host above those thresholds,
every $t\ge2$ satisfying \texttt{Big}$(t)$, and every $B\ge0$
for which
\[
 \mathtt{LayerLe}(\chi,\ell,j(\ell),t,B)
                      \quad(\ell\in\mathcal B_\chi)
 \tag{HV.42}
\]
holds, every original CM field $E$ of that host satisfies
\[
                       \mathtt{VVplusLe}
                               \bigl(E,t,(Nt)^C+CB\bigr).
 \tag{HV.43}
\]
With TW8, OB8 and the already reconstructed VW assembly,
the same four-count premise therefore implies the full
cohomological window conclusion
\[
 \forall C_t\ge1\ \forall\text{ primitive }\psi,
                         \mathtt{VWinAt}(\psi,C_t).
 \tag{HV.44}
\]
The previously stated prime-pair inputs of VW retain their
stated role.  In particular its $C_t=401$ specialization
retains the fixed pair and the cutoff exponent $400$.

\begin{proof}
By HV8 and HV9 the four-count premise supplies \texttt{SumIn}
and \texttt{SumOut}; HV5, HV6 and CZ6 have already supplied
their other inputs.  Fix $\chi$ and the whole function $j$.
Choose $C_{\rm out},r_1,r_2$ from HV9 and, for each prime
of the finite set $\mathcal B_\chi$, choose the constant
$C_{\rm in}(\ell,j(\ell))$ from HV8.  Put
\[
 D=\sum_{\ell\in\mathcal B_\chi}C_{\rm in}(\ell,j(\ell)),
 \qquad k=|\mathcal B_\chi|,
 \qquad C=C_{\rm out}+D+k+1.
 \tag{HV.45}
\]
These are choices before the host, $t,B,E,\Omega$ and local
data.  For any subsequent choices satisfying (HV.42), put
$M=Nt\ge2$ and split an arbitrary finite set $S$ of odd primes
into its intersection with $\mathcal B_\chi$ and its complement.
The sum in (HV.43) over $S$ is bounded by
\[
 \begin{split}
 M^{C_{\rm out}}+
 \sum_{\ell\in S\cap\mathcal B_\chi}
       \bigl(M^{C_{\rm in}(\ell,j(\ell))}
                          +C_{\rm in}(\ell,j(\ell))B\bigr)
 &\le (k+1)M^{C_{\rm out}+D}+DB\\
 &\le M^C+CB.
 \end{split}
 \tag{HV.46}
\]
Here each $C_{\rm in}$ is positive, $B\ge0$, and
$k+1\le2^k\le M^k$ (also valid at $k=0$).
The two component results prove all individual finiteness
assertions.  The field and local data were arbitrary,
and (HV.46) is uniform over every finite $S$.
This is exactly \texttt{VVplusLe}, hence the literal
\texttt{PointHV}, with the full quantifier order in (HV.43).

For (HV.44), use the single constant record and universal
\texttt{ThmBab} supplied by TW8, and the odd-minus
\texttt{ThmBcFree} supplied by OB8.  Substitution of (HV.41)
into VW3--VW4 now fills their remaining odd-plus premise.
Those proofs already keep the same original field, all the
prescribed layer models, one common window and the fixed
prime pair.  Their rescaling with $C_t=401$ is unchanged.
This proves the asserted immediate propagation.
\end{proof}

\paragraph{Scope at the end of the chain.}
HV proves the two auxiliary counts and every step after the
four-count cut.  In particular the separate
\texttt{PointHV}-assembly task is complete, including bad
degree primes and the coefficient and field comparisons.
The full human reconstruction still needs the actual
\texttt{CountI}, \texttt{CountII}, \texttt{CountIII} and
\texttt{CountV}, or complete sufficient suppliers for them.
The common arithmetic family, genuine theta/local installation,
denominator tables, arithmetic gap and exact metric/spectral
comparisons remain the independent global obligations already
identified in the audit.  Neither (HV.41) nor (HV.44) asserts
that those inputs have been constructed.
